Um belo faz-de-conta;
    é tudo que a ciência pode ser:
        mas a busca dos \textit{humanos} pela \textbf{Verdade} tem-lhe subtraído tamanha potência.

Para falar de \textbf{Verdade}, é preciso tratá-la como poesia;
    não se pode buscá-la, apenas a viver atentando-se a seus movimentos fugidios:
        o tal "mundo real" que se almeja descrever com ela é um imenso e inconstante mar de ondas, existências e mistérios.

Em suas águas não se encontram linhas, fronteiras, limites;
    afinal, quem poderia dizer onde começa e termina uma onda?:
        a gente pode até olhar, observar a imensidão acachapante do oceano, e enxergar distintas ondas por sua superfície, mas traçar seus contornos é uma tarefa cuja própria imaginação já me traz certo desconcerto ao pensamento.

A questão evolui no dar-se-conta de que efetivamente identificamos uma onda em meio ao mar, ou a outras ondas, ou existências;
    o problema é que não há identidade sem margens que a separem do que não está sendo identificado:
        se digo que \textbf{\textit{a surfista do Brasil pegou uma onda fraca em sua última bateria do torneio}}, identifico a onda como algo que não é a surfista, nem o Brasil, nem a bateria, nem o torneio --- separo as coisas através das palavras.

Nomear é uma operação de estabelecimento de fronteiras-identidades, assim como categorizar o é de fronteiras-tipos;
    nomes e categorias são mapas que identificam -- separam -- e tipificam -- reagrupam -- ondas no mar da realidade:
        a onda fraca da frase acima é única entre as existências invocadas, está separada das demais --- mas ela é fraca, isto é, entre as ondas fracas ela encontra sua tribo.

Me parece haver aqui \textit{um tabuleiro comum a todos os jogos};
    penso em comunicação como um processo de composição de comunalidade:
        com a frase da surfista, componho-com as existências alfabetizadas em língua portuguesa um \textbf{mapa-em-comum}, comunico-com elas.

O tabuleiro comum da comunicação estabelece a separação e o reagrupamento como componentes ancestrais de todas as jogadas possíveis, realizadas ou não;
    todos os jogos são jogos de diferenciação-e-repetição:
        "onda fraca" destaca de um acontecimento específico -- circunstancialmente hipotético -- -com o qual a frase comunica, um de seus elementos --- e o reagrupa junto a elementos de outros acontecimentos, faz com que eles se repitam numa cena única.

O jogo da ciência moderna compõe-com o jogo-ancestral da diferenciação-e-repetição regras -- epistemologias, modos de saber e de afirmar, isto é, de jogar -- proliferantes de repetições;
    o específico diferenciado participa das jogadas como algo a ser consumido e obliterado a fim de se compor-com um conjunto de acontecimentos e elementos únicos um conjunto de categorias que diminuam a imprevisibilidade -- a unicidade circunstancial -- dos acontecimentos por vir:
        trata-se de uma estratégia de composição do comum orientada e mediada pelo apagamento e superação do diferente.

O objetivo dessa brincadeira é alcançar a \textbf{Verdade}, a cartografia de todo o mar, como se fosse possível delinear exatamente onde uma de suas ondas começa e outra termina;
    um jogo de faz-de-conta no qual já se pôde encontrar a beleza, ao menos, de sinérgicas repetições, de um aumento progressivo dos avistamentos de baleias a saltarem sobre a superfície das águas:
        hoje, contudo, elas se tornaram tantas que mal se consegue ver o mar e suas ondas --- a proliferação desses bichos justapondo-se em nossos avistamentos tem alimentado uma cacofonia silenciosa que progressivamente diminui a capacidade dos jogadores de identificar o contorno, a posição e a trajetória de cada baleia.

Em sua ânsia pelo progresso, o homem moderno construiu por conta própria o caminho que o levou de volta ao seu ponto de partida;
    sua temporalidade linear revelou-se a ação de um trickster a pregar-lhe uma peça:
        "siga em frente toda a vida", ele lhes disse, entregando-lhes a linha reta --- uma armadilha ardilosa que nunca falha em se revelar um bumerangue.

No mundo que almejo descrever contam-se piadas sobre o terraplanismo temporal dos modernos;
    a linearidade do tempo é uma ilusão de ótica da mesma ordem que a planitude da Terra:
        a orientação pelo progresso tem constituído um novo mar, de baleias, sobre as águas marinhas --- marchando sempre em linha reta, os soldados da modernidade progridem em tropicões cada vez mais frequentes com a temporalidade espiralar do chão em que pisam.

O jogo dos modernos vai se esgotando, e um novo pacto com o jogo-ancestral começa a se compor;
    variadas linhas retas já se tem cruzado revelando movimentos mais complexos-coletivos-simpoiéticos na multidão de baleias:
        proponho que este modo \textit{outro} de jogar articule os avistamentos de existências saltitantes aos de existências imanentes como as ondas.

A composição do \textbf{\textit{comum}} -- isto é, a comunicação -- é a articulação do \textbf{diferente} com o \textit{igual};
    \textbf{\textit{palavras}} vestem \textbf{coisas únicas em si} como \textit{coisas semelhantes ao outro}:
